\appendix

\section{Proof Details}\label{sec:proof-details}

\subsection{Root Feasibility and Finite Pivot Conditioning}
\label{app:step-001}

\begin{lemma}[Supported-root restriction and compact feasible set]
\label{lem:step-001-root-feasible}
Under Assumption~\ref{assump:shared-pfaffian-chain} and the basic conditions
$N\geq1$, $R>0$, and compact $\Theta$, the set
\[
K_R=\{\theta\in\Theta:|b(\theta)|\leq R\|F(\theta)\|_1\}
\]
is compact.  For every $\alpha\in[-R,R]^N$ and $\theta\in\Theta$,
$\phi_\alpha(\theta)=0$ implies $\theta\in K_R$.  Consequently, if
$K_R=\varnothing$, no coefficient vector in $[-R,R]^N$ has a root in
$\Theta$, and $\Gamma_{\rm piv}(b,F;R)=0$.
\end{lemma}

\begin{proof}
Assumption~\ref{assump:shared-pfaffian-chain} makes
$b,F_1,\ldots,F_N$ continuous.  Hence
\[
\theta\longmapsto
|b(\theta)|-R\sum_{i=1}^N|F_i(\theta)|
\]
is continuous, and $K_R$ is its inverse image of the closed set
$(-\infty,0]$.  Thus $K_R$ is closed in compact $\Theta$ and is compact.

Fix $\alpha\in[-R,R]^N$ and suppose $\phi_\alpha(\theta)=0$.  Then
\[
\begin{aligned}
|b(\theta)|
&=\left|\sum_{i=1}^N\alpha_iF_i(\theta)\right|\\
&\leq\sum_{i=1}^N|\alpha_i|\,|F_i(\theta)|\\
&\leq R\sum_{i=1}^N|F_i(\theta)|
=R\|F(\theta)\|_1.
\end{aligned}
\]
Therefore $\theta\in K_R$.  If $K_R$ is empty, this implication excludes
every supported root.  The empty-set branch in the definition of
$\Gamma_{\rm piv}$ then gives $\Gamma_{\rm piv}(b,F;R)=0$.
\end{proof}

\begin{lemma}[Compact available-pivot margin]
\label{lem:step-001-pivot-margin}
Under Assumptions~\ref{assump:shared-pfaffian-chain} and
\ref{assump:no-forced-root} and
Lemma~\ref{lem:step-001-root-feasible}, if $K_R\neq\varnothing$, then
$F(\theta)\neq0$ for every $\theta\in K_R$, and
\[
\rho:=\min_{\theta\in K_R}\max_{1\leq j\leq N}|F_j(\theta)|
\]
is well-defined and strictly positive.  Consequently, for every
$\theta\in K_R$, at least one coordinate $j$ satisfies
$|F_j(\theta)|\geq\rho>0$.
\end{lemma}

\begin{proof}
Take $\theta\in K_R$.  If $F(\theta)=0$, then feasibility yields
\[
|b(\theta)|\leq R\|F(\theta)\|_1=0,
\]
so $b(\theta)=0$.  This contradicts
Assumption~\ref{assump:no-forced-root}.  Thus $F(\theta)\neq0$ throughout
$K_R$, and
\[
\max_{1\leq j\leq N}|F_j(\theta)|>0
\quad\text{for every }\theta\in K_R.
\]

The maximum of finitely many continuous coordinate magnitudes is continuous.
By Lemma~\ref{lem:step-001-root-feasible}, $K_R$ is nonempty and compact, so
this function attains its minimum there.  If the attained minimum were zero,
its minimizer would have $F_j=0$ for every $j$, contradicting the
nonvanishing just proved.  Hence $\rho>0$.  At each point, one of the
finitely many coordinates attains the maximum, proving the final assertion.
\end{proof}

\begin{proposition}[Fixed-family pivot-variation bound]
\label{prop:step-001-fixed-family-bound}
Under Assumption~\ref{assump:shared-pfaffian-chain} and
Lemma~\ref{lem:step-001-pivot-margin}, suppose $K_R\neq\varnothing$, and
define
\[
B_0:=\max_{\theta\in\Theta}
\max\{|b(\theta)|,|F_1(\theta)|,\ldots,|F_N(\theta)|\},
\]
\[
B_1:=\max_{\theta\in\Theta}
\max\{|b'(\theta)|,|F_1'(\theta)|,\ldots,|F_N'(\theta)|\}.
\]
Then $B_0,B_1<\infty$, and for every $\theta\in K_R$,
\[
\min_{1\leq k\leq N}V_k(\theta)
\leq
\frac{2B_0B_1}{\rho^2}\bigl(1+R(N-1)\bigr).
\]
Consequently,
\[
\Gamma_{\rm piv}(b,F;R)
\leq
\frac{2B_0B_1}{\rho^2}\bigl(1+R(N-1)\bigr)<\infty.
\]
\end{proposition}

\begin{proof}
Assumption~\ref{assump:shared-pfaffian-chain} makes each of
$b,F_i,b',F_i'$ continuous on compact $\Theta$.  The extreme-value theorem
therefore makes both displayed maxima finite.

Fix $\theta\in K_R$.  By Lemma~\ref{lem:step-001-pivot-margin}, choose a
coordinate $j$ attaining $\max_k|F_k(\theta)|$.  Then
$|F_j(\theta)|\geq\rho>0$, so $V_j(\theta)$ is evaluated on its finite-pivot
branch.  For any $g\in\{b,F_i:i\neq j\}$, the quotient rule gives
\[
\left(\frac g{F_j}\right)'(\theta)
=\frac{g'(\theta)F_j(\theta)-g(\theta)F_j'(\theta)}{F_j(\theta)^2},
\]
and therefore
\[
\begin{aligned}
\left|\left(\frac g{F_j}\right)'(\theta)\right|
&=\frac{|g'(\theta)F_j(\theta)-g(\theta)F_j'(\theta)|}
{|F_j(\theta)|^2}\\
&\leq
\frac{|g'(\theta)|\,|F_j(\theta)|
      +|g(\theta)|\,|F_j'(\theta)|}{|F_j(\theta)|^2}\\
&\leq\frac{B_1B_0+B_0B_1}{\rho^2}
=\frac{2B_0B_1}{\rho^2}.
\end{aligned}
\]
Applying this bound once with $g=b$ and once for each of the $N-1$ indices
$i\neq j$ yields
\[
\begin{aligned}
V_j(\theta)
&=\left|\left(\frac b{F_j}\right)'(\theta)\right|
+R\sum_{i\neq j}
 \left|\left(\frac{F_i}{F_j}\right)'(\theta)\right|\\
&\leq\frac{2B_0B_1}{\rho^2}\bigl(1+R(N-1)\bigr).
\end{aligned}
\]
The minimum over all pivots is at most this particular finite value.  Taking
the supremum over $K_R$ proves the asserted bound.  When $N=1$, the ratio
sum is empty and the same derivation gives the factor
$1+R(N-1)=1$.

The estimate is for the fixed deterministic family.  It does not estimate
$B_0$, $B_1$, or $\rho^{-1}$ from
$(q,M,\Delta_{\rm rnd},\Delta_{\rm aff})$ or other Pfaffian-format data.
\end{proof}

\begin{proof}[Completion of the feasibility and finiteness argument]
Lemma~\ref{lem:step-001-root-feasible} proves that $K_R$ is compact and
contains every root produced by a coefficient in $[-R,R]^N$.  If
$K_R=\varnothing$, it also proves that the supported-root event is empty,
and the definition gives $\Gamma_{\rm piv}=0$.

If $K_R\neq\varnothing$, Lemma~\ref{lem:step-001-pivot-margin} derives the
strict margin
$\rho=\min_{K_R}\max_j|F_j|>0$ from
Assumption~\ref{assump:no-forced-root}.  Thus individual coordinates may
vanish, but at least one is uniformly available at every root-feasible
point.  Proposition~\ref{prop:step-001-fixed-family-bound} then gives
\[
\min_jV_j(\theta)
\leq2B_0B_1\rho^{-2}\bigl(1+R(N-1)\bigr),
\qquad \theta\in K_R,
\]
and hence $\Gamma_{\rm piv}(b,F;R)<\infty$.  The two cases prove root
feasibility, qualitative pivot availability, and fixed-family finiteness,
without asserting polynomial control by Pfaffian-format descriptors.
\end{proof}

\subsection{Measurable Adaptive Pivot Charts}
\label{app:step-002}

\begin{lemma}[Borel extended pivot velocities]
\label{lem:step-002-borel-velocities}
Under Assumption~\ref{assump:shared-pfaffian-chain}, for each
$j\in\{1,\ldots,N\}$, the extension $V_j:\Theta\to[0,+\infty]$, equal to
\[
\left|\left(\frac b{F_j}\right)'\right|
+R\sum_{i\neq j}\left|\left(\frac{F_i}{F_j}\right)'\right|
\quad\text{on }U_j=\{F_j\neq0\},
\]
and equal to $+\infty$ on $\Theta\setminus U_j$, is Borel.  Moreover, for
all $j,k$, the sets $\{V_j<V_k\}$ and $\{V_j\leq V_k\}$ are Borel subsets
of $\Theta$.
\end{lemma}

\begin{proof}
Assumption~\ref{assump:shared-pfaffian-chain} makes $F_j$ continuous, so
$U_j$ is open in $\Theta$.  Since the compact interval $\Theta$ is Borel in
$\mathbb R$, every relative Borel subset used below is also Borel in the
ambient line.  For $g=b$, and separately for every $g=F_i$ with $i\neq j$,
the quotient rule gives on $U_j$
\[
\left(\frac g{F_j}\right)'
=\frac{g'F_j-gF_j'}{F_j^2}.
\]
The numerator and denominator are continuous there, and the denominator does
not vanish.  Every displayed quotient derivative is therefore continuous on
$U_j$, as is their finite absolute-value sum defining $V_j$.

For each finite $a\in\mathbb R$,
\[
\{\theta\in\Theta:V_j(\theta)<a\}
=U_j\cap
\left\{\theta\in U_j:
\left|\left(\frac b{F_j}\right)'(\theta)\right|
+R\sum_{i\neq j}
\left|\left(\frac{F_i}{F_j}\right)'(\theta)\right|<a
\right\}.
\]
The right-hand side is Borel in $\Theta$, and
$\{V_j=+\infty\}=\Theta\setminus U_j$ is closed.  Thus $V_j$ is Borel as
an extended-real map.

Strict comparison can be checked at a rational level:
\[
\{V_j<V_k\}
=\bigcup_{r\in\mathbb Q}
\bigl(\{V_j<r\}\cap\{V_k>r\}\bigr).
\]
If $V_j(\theta)<V_k(\theta)$, then $V_j(\theta)$ is finite and density of
$\mathbb Q$ supplies a finite separator, including when
$V_k(\theta)=+\infty$; the reverse implication is immediate.  Hence the
strict comparison is Borel.  Finally,
\[
\{V_j\leq V_k\}=\Theta\setminus\{V_k<V_j\}
\]
is Borel as well.
\end{proof}

\begin{proposition}[Measurable least-pivot partition]
\label{prop:step-002-selector-partition}
Under Assumptions~\ref{assump:shared-pfaffian-chain} and
\ref{assump:no-forced-root}, Lemmas~\ref{lem:step-001-root-feasible} and
\ref{lem:step-001-pivot-margin},
Proposition~\ref{prop:step-001-fixed-family-bound}, and
Lemma~\ref{lem:step-002-borel-velocities}, define for every
$\theta\in K_R$
\[
j_*(\theta)=\min\operatorname*{argmin}_{1\leq k\leq N}V_k(\theta),
\qquad
E_j=\{\theta\in K_R:j_*(\theta)=j\}.
\]
Then $j_*:K_R\to\{1,\ldots,N\}$, with the discrete codomain, is Borel, and
$(E_j)_{j=1}^N$ is a disjoint Borel partition of $K_R$.  Moreover, for every
$\theta\in E_j$,
\[
F_j(\theta)\neq0,
\qquad
V_j(\theta)=\min_{1\leq k\leq N}V_k(\theta)
\leq\Gamma_{\rm piv}(b,F;R)<\infty.
\]
When $K_R=\varnothing$, the selector is the empty map, all cells are empty,
and the same conclusions hold vacuously.
\end{proposition}

\begin{proof}
Suppose first that $K_R\neq\varnothing$.  Lemma
\ref{lem:step-001-pivot-margin} and
Proposition~\ref{prop:step-001-fixed-family-bound} ensure that at each
$\theta\in K_R$ at least one $V_k(\theta)$ is finite.  A finite family of
extended-real values therefore has a finite minimum and a unique least index
attaining it, so $j_*$ is well-defined.

The least-index rule has the exact cell representation
\begin{equation}
E_j
=K_R\cap\bigcap_{k<j}\{V_j<V_k\}
       \cap\bigcap_{k>j}\{V_j\leq V_k\}.
\label{eq:least-pivot-cell}
\end{equation}
A minimizing index $j$ is least exactly when its value is strictly below
every earlier-index value and no larger than every later-index value.
Lemma~\ref{lem:step-001-root-feasible} makes $K_R$ Borel, and
Lemma~\ref{lem:step-002-borel-velocities} makes every comparison set in
\eqref{eq:least-pivot-cell} Borel.  Thus each $E_j$ is Borel.  The unique
least minimizer places every point of $K_R$ in exactly one cell, so the cells
are pairwise disjoint and their union is $K_R$.  Since
$j_*^{-1}(\{j\})=E_j$ for every singleton in the finite discrete codomain,
the selector is Borel.

If $\theta\in E_j$, then $V_j(\theta)$ equals the finite minimum.  By the
definition of the extension, $F_j(\theta)=0$ would force
$V_j(\theta)=+\infty$, a contradiction.  Hence the selected pivot is
nonzero.  The definition of $\Gamma_{\rm piv}$ on nonempty $K_R$ gives
\[
V_j(\theta)=\min_kV_k(\theta)
\leq\sup_{\vartheta\in K_R}\min_kV_k(\vartheta)
=\Gamma_{\rm piv}(b,F;R),
\]
and Proposition~\ref{prop:step-001-fixed-family-bound} proves finiteness.

If $K_R=\varnothing$, the unique map from the empty domain is Borel, every
cell is empty, and all pointwise statements are vacuous.  For $N=1$, both
intersections in \eqref{eq:least-pivot-cell} are empty intersections, so
$E_1=K_R$; the preceding argument still gives $F_1\neq0$ on nonempty
$K_R$.
\end{proof}

\begin{lemma}[Finite-pivot chart exhaustion]
\label{lem:step-002-chart-exhaustion}
Under Assumption~\ref{assump:shared-pfaffian-chain} and
Proposition~\ref{prop:step-002-selector-partition}, for every
$j\in\{1,\ldots,N\}$ and integer $m\geq1$,
\[
E_{j,m}=E_j\cap\{\theta\in\Theta:|F_j(\theta)|\geq1/m\}
\]
is Borel, $E_{j,m}\subseteq E_{j,m+1}$, and
\[
\bigcup_{m=1}^{\infty}E_{j,m}=E_j.
\]
\end{lemma}

\begin{proof}
The cell $E_j$ is Borel by
Proposition~\ref{prop:step-002-selector-partition}.  Continuity of $F_j$
makes $\{|F_j|\geq1/m\}$ closed in $\Theta$, so $E_{j,m}$ is Borel.  Since
$1/(m+1)\leq1/m$,
\[
\{|F_j|\geq1/m\}\subseteq\{|F_j|\geq1/(m+1)\},
\]
which proves monotonicity.

Every $E_{j,m}$ is contained in $E_j$.  Conversely, if $\theta\in E_j$,
Proposition~\ref{prop:step-002-selector-partition} gives
$|F_j(\theta)|>0$.  Choose an integer
$m\geq\max\{1,|F_j(\theta)|^{-1}\}$.  Then
$1/m\leq|F_j(\theta)|$, so $\theta\in E_{j,m}$.  This proves the union
identity, including arbitrarily small selected pivots and empty cells.
\end{proof}

\begin{proposition}[Exact affine pivot chart and material velocity]
\label{prop:step-002-exact-chart}
Under Assumption~\ref{assump:shared-pfaffian-chain} and
Proposition~\ref{prop:step-002-selector-partition}, fix
$j\in\{1,\ldots,N\}$, $\theta\in E_j$, and
$\beta=(\beta_i)_{i\neq j}\in[-R,R]^{N-1}$.  Then
\[
T_j(\theta,\beta)
=-\frac{b(\theta)+\sum_{i\neq j}\beta_iF_i(\theta)}{F_j(\theta)}
\]
is well-defined.  It is the unique value of the original coefficient
$\alpha_j$ for which the original vector with $\alpha_{-j}=\beta$ satisfies
$\phi_\alpha(\theta)=0$.  The chart is the restriction of a $C^1$, hence
Borel, function on $U_j\times\mathbb R^{N-1}$, and its material partial
derivative, with $\beta$ fixed, satisfies
\begin{equation}
\partial_\theta T_j(\theta,\beta)
=-\left(\frac b{F_j}\right)'(\theta)
-\sum_{i\neq j}\beta_i
 \left(\frac{F_i}{F_j}\right)'(\theta),
\label{eq:chart-derivative}
\end{equation}
and
\begin{equation}
|\partial_\theta T_j(\theta,\beta)|
\leq V_j(\theta)\leq\Gamma_{\rm piv}(b,F;R).
\label{eq:chart-speed}
\end{equation}
For $N=1$, $\beta$ is the unique empty tuple, both sums are empty, and the
statement reads $T_1=-b/F_1$ and
$|\partial_\theta T_1|=V_1\leq\Gamma_{\rm piv}$.
\end{proposition}

\begin{proof}
Proposition~\ref{prop:step-002-selector-partition} gives
$F_j(\theta)\neq0$, so the quotient is well-defined.  If
$\alpha_{-j}=\beta$ and $\alpha_j=T_j(\theta,\beta)$, then in the original
coefficient coordinates
\begin{equation}
\begin{aligned}
\phi_\alpha(\theta)
&=b(\theta)+\sum_{i\neq j}\beta_iF_i(\theta)
  +T_j(\theta,\beta)F_j(\theta)\\
&=b(\theta)+\sum_{i\neq j}\beta_iF_i(\theta)
 -\left(b(\theta)+\sum_{i\neq j}\beta_iF_i(\theta)\right)\\
&=0.
\end{aligned}
\label{eq:chart-root-identity}
\end{equation}
Conversely, if $\phi_\alpha(\theta)=0$ and $\alpha_{-j}=\beta$, division by
the same nonzero $F_j(\theta)$ gives
$\alpha_j=T_j(\theta,\beta)$.  Thus
\eqref{eq:chart-root-identity} is exact in the original coordinates, and
uniqueness follows from the nonzero coefficient of $\alpha_j$.

On $U_j\times\mathbb R^{N-1}$, rewrite
\[
T_j(\theta,\beta)
=-\frac{b(\theta)}{F_j(\theta)}
-\sum_{i\neq j}\beta_i\frac{F_i(\theta)}{F_j(\theta)}.
\]
All ratios are $C^1$ there.  Holding $\beta$ fixed and applying the quotient
rule gives \eqref{eq:chart-derivative}, with no derivative of $\beta$.  The
triangle inequality and $|\beta_i|\leq R$ give
\[
\begin{aligned}
|\partial_\theta T_j(\theta,\beta)|
&\leq
\left|\left(\frac b{F_j}\right)'(\theta)\right|
+\sum_{i\neq j}|\beta_i|
 \left|\left(\frac{F_i}{F_j}\right)'(\theta)\right|\\
&\leq
\left|\left(\frac b{F_j}\right)'(\theta)\right|
+R\sum_{i\neq j}
 \left|\left(\frac{F_i}{F_j}\right)'(\theta)\right|\\
&=V_j(\theta).
\end{aligned}
\]
The final inequality in \eqref{eq:chart-speed} is the selected-cell
conclusion of Proposition~\ref{prop:step-002-selector-partition}.

When $N=1$, the remaining-coordinate space is $\mathbb R^0$, the cube
$[-R,R]^0$ consists of the empty tuple and has zero-dimensional volume one,
and every sum over $i\neq1$ is zero.  The same algebra gives
$T_1=-b/F_1$, \eqref{eq:chart-root-identity}, and
\[
|\partial_\theta T_1|
=\left|\left(\frac b{F_1}\right)'\right|
=V_1\leq\Gamma_{\rm piv}.
\]
\end{proof}

\begin{proof}[Completion of the measurable-chart argument]
Lemma~\ref{lem:step-002-borel-velocities} supplies every comparison needed
for finite minimization.  Proposition~\ref{prop:step-002-selector-partition}
uses the finite-pivot conclusions of
Lemmas~\ref{lem:step-001-root-feasible} and
\ref{lem:step-001-pivot-margin} and
Proposition~\ref{prop:step-001-fixed-family-bound} to implement least-index
ties by \eqref{eq:least-pivot-cell}.  It yields a disjoint Borel partition
with $F_j\neq0$ and
$V_j=\min_kV_k\leq\Gamma_{\rm piv}$ on $E_j$.

Lemma~\ref{lem:step-002-chart-exhaustion} gives the literal nested identity
$E_{j,m}\uparrow E_j$; the index $m$ only exposes finite denominators.
Proposition~\ref{prop:step-002-exact-chart} inserts the unique original
coefficient solving the original equation and proves
\eqref{eq:chart-derivative}--\eqref{eq:chart-speed}, including the empty
tuple when $N=1$.  The empty-$K_R$ branch has empty cells.  These results
provide the measurable graph and Jacobian interface used below, without a
globally prescribed coordinate or an $N$-fold chart loss.
\end{proof}

\subsection{Measurable Root-Event Volume by Adaptive Charts}
\label{app:step-003}

\begin{lemma}[Measurable exhausted coefficient charts]
\label{lem:step-003-measurable-charts}
Under Assumption~\ref{assump:shared-pfaffian-chain},
Proposition~\ref{prop:step-002-selector-partition},
Lemma~\ref{lem:step-002-chart-exhaustion}, and
Proposition~\ref{prop:step-002-exact-chart}, fix an interval
$I\subseteq\Theta$.  For $j\in\{1,\ldots,N\}$, order the remaining
coordinates as
\[
\beta=(\beta_1,\ldots,\beta_{j-1},\beta_{j+1},\ldots,\beta_N)
\in B_j:=[-R,R]^{N-1},
\]
and define, in original coefficient order,
\begin{equation}
\Psi_j(\theta,\beta)
=
(\beta_1,\ldots,\beta_{j-1},
T_j(\theta,\beta),
\beta_{j+1},\ldots,\beta_N).
\label{eq:step-003-insertion-map}
\end{equation}
For every integer $m\geq1$, set
\begin{equation}
D_{j,m}
=\{(\theta,\beta):
\theta\in I\cap E_{j,m},\ \beta\in B_j,\
|T_j(\theta,\beta)|\leq R\}.
\label{eq:step-003-exhausted-domain}
\end{equation}
Then $D_{j,m}$ is Borel in $\mathbb R^N$, $\Psi_j$ is the restriction of a
$C^1$, hence locally Lipschitz, map on an open neighborhood of every point
of $D_{j,m}$, and $\Psi_j(D_{j,m})$ is analytic and Lebesgue measurable.
The event $S_I$ is analytic and Lebesgue measurable.  These assertions
include interval endpoints, the closed coefficient-cube boundary, and the
$N=1$ convention $B_1=[-R,R]^0=\{()\}$.
\end{lemma}

\begin{proof}
Write $\Theta=[\theta_-,\theta_+]$.  A $C^1$ function $g$ on this compact
interval has the explicit $C^1$ extension
\begin{equation}
\widetilde g(x)=
\begin{cases}
g(\theta_-)+g'(\theta_-)(x-\theta_-),&x<\theta_-,\\
g(x),&\theta_-\leq x\leq\theta_+,\\
g(\theta_+)+g'(\theta_+)(x-\theta_+),&x>\theta_+.
\end{cases}
\label{eq:step-003-c1-extension}
\end{equation}
The values and first derivatives match at both endpoints.  If
$\theta_-=\theta_+$, use the single affine expression
$g(\theta_-)+g'(\theta_-)(x-\theta_-)$.  Apply
\eqref{eq:step-003-c1-extension} separately to $b$ and to
$F_1,\ldots,F_N$.  On $\Theta$, these extensions preserve the original
functions and their derivatives exactly.

For each $j$, define the open set
\[
O_j=\{(\theta,\beta)\in\mathbb R\times\mathbb R^{N-1}:
\widetilde F_j(\theta)\neq0\}.
\]
The quotient formula for $T_j$, now using the extended functions, defines a
$C^1$ function on $O_j$, and
\eqref{eq:step-003-insertion-map} defines a $C^1$, locally Lipschitz map
there.  This requires no lower bound for one pivot over all of $E_j$:
Proposition~\ref{prop:step-002-selector-partition} gives the pointwise fact
$F_j\neq0$ on $E_j$, while $E_{j,m}$ records $|F_j|\geq1/m$.

Lemma~\ref{lem:step-002-chart-exhaustion} makes $E_{j,m}$ Borel.
Intervals and the closed cube $B_j$ are Borel.  On
$E_{j,m}\times B_j\subset O_j$, the condition $|T_j|\leq R$ is the inverse
image of a closed interval under a continuous function.  Hence
$D_{j,m}$ is Borel.  The restriction of $\Psi_j$ to this Borel set is
Borel, so its image is analytic and therefore Lebesgue measurable; here we
use the Borel-image and universal-measurability theorems for analytic sets
in Polish spaces \citep{kechris1995}.

The target event is measurable independently of the coverage identity proved
below.  Indeed,
\[
Z_I=\{(\alpha,\theta)\in[-R,R]^N\times I:
b(\theta)+\langle\alpha,F(\theta)\rangle=0\}
\]
is Borel because affine evaluation is continuous on the Borel product
restriction.  Its coefficient projection is exactly $S_I$.  A coordinate
projection of a Borel subset of a product of Euclidean spaces is analytic,
so $S_I$ is analytic and Lebesgue measurable \citep{kechris1995}.

No boundary was removed: \eqref{eq:step-003-exhausted-domain} uses the actual
interval $I$, the closed beta cube, and $|T_j|\leq R$, including
$T_j=\pm R$.  The extension \eqref{eq:step-003-c1-extension} supplies an
ambient neighborhood even at endpoints of $\Theta$.  When $N=1$,
$\mathbb R^0$ is a one-point Euclidean space, $B_1$ is that point, and all
claims reduce to their one-dimensional versions.
\end{proof}

\begin{proposition}[Noninjective area bound in original coordinates]
\label{prop:step-003-area-bound}
Under Assumption~\ref{assump:shared-pfaffian-chain},
Proposition~\ref{prop:step-002-exact-chart}, and
Lemma~\ref{lem:step-003-measurable-charts}, use domain-coordinate order
\[
(\theta,\beta_1,\ldots,\beta_{j-1},\beta_{j+1},\ldots,\beta_N)
\]
and target-coordinate order $(\alpha_1,\ldots,\alpha_N)$.  Then
\begin{equation}
\det D\Psi_j(\theta,\beta)
=(-1)^{j-1}\partial_\theta T_j(\theta,\beta),
\qquad
J_N\Psi_j=|\partial_\theta T_j|.
\label{eq:step-003-jacobian}
\end{equation}
For every $j$ and $m\geq1$,
\begin{equation}
\begin{aligned}
\lambda_N(\Psi_j(D_{j,m}))
&\leq
\int_{D_{j,m}}|\partial_\theta T_j(\theta,\beta)|\,d\theta\,d\beta\\
&\leq
(2R)^{N-1}\int_{I\cap E_{j,m}}V_j(\theta)\,d\theta.
\end{aligned}
\label{eq:step-003-area-bound}
\end{equation}
The first inequality permits finite or infinite multiplicity and requires no
simple-root, transversality, or injectivity premise.  For $N=1$,
\eqref{eq:step-003-jacobian} is the ordinary derivative and
\eqref{eq:step-003-area-bound} uses
$\lambda_0(B_1)=1=(2R)^0$.
\end{proposition}

\begin{proof}
For $i\neq j$, the $\alpha_i$ component of
\eqref{eq:step-003-insertion-map} is $\beta_i$, while the $\alpha_j$
component is $T_j$.  Let $P_j$ move target row $j$ to the first row and
preserve the order of all other rows.  It performs $j-1$ transpositions, so
$\det P_j=(-1)^{j-1}$.  Direct differentiation gives
\[
P_jD\Psi_j
=
\begin{pmatrix}
\partial_\theta T_j&\nabla_\beta T_j\\
0&I_{N-1}
\end{pmatrix}.
\]
Consequently,
\[
\det(P_jD\Psi_j)=\partial_\theta T_j,
\qquad
\det D\Psi_j=(-1)^{j-1}\partial_\theta T_j.
\]
This proves \eqref{eq:step-003-jacobian}, including the permutation sign.
The entries $\partial_{\beta_i}T_j=-F_i/F_j$ occupy the upper-right block
and do not alter the determinant.

We next verify the measurable restriction needed by the area formula.  For
fixed $j,m$,
\[
C_{j,m}
=\{\theta\in\Theta:|F_j(\theta)|\geq1/m\}\times B_j
\]
is compact, lies in $O_j$, and contains $D_{j,m}$.  Finitely many open balls
compactly contained in $O_j$, on each of which the $C^1$ map $\Psi_j$ is
Lipschitz, cover $C_{j,m}$.  Intersect the balls with $D_{j,m}$ and
successively remove earlier balls.  This produces a finite disjoint Borel
partition $(A_\ell)_\ell$ of $D_{j,m}$, with $\Psi_j$ Lipschitz on each
piece.  The equal-dimensional Euclidean area formula in multiplicity form
\citep{azaisWschebor2009}, image subadditivity, and domain disjointness give
\begin{equation}
\begin{aligned}
\lambda_N(\Psi_j(D_{j,m}))
&\leq\sum_\ell\lambda_N(\Psi_j(A_\ell))\\
&\leq\sum_\ell\int_{A_\ell}J_N\Psi_j
=\int_{D_{j,m}}J_N\Psi_j.
\end{aligned}
\label{eq:step-003-localized-area}
\end{equation}
Thus localization contributes no numerical factor.  Equivalently, one may
exhaust an arbitrary ambient open set by countably many compact
localizations and apply the same disjoint refinement.  The threshold $1/m$
is only the denominator exhaustion level.

By \eqref{eq:step-003-jacobian},
Proposition~\ref{prop:step-002-exact-chart}, and Tonelli's theorem,
\begin{equation}
\begin{aligned}
\int_{D_{j,m}}J_N\Psi_j\,d\theta\,d\beta
&=\int_{I\cap E_{j,m}}\int_{B_j}
{\bf1}_{\{|T_j(\theta,\beta)|\leq R\}}
|\partial_\theta T_j(\theta,\beta)|\,d\beta\,d\theta\\
&\leq\int_{I\cap E_{j,m}}\int_{B_j}
V_j(\theta)\,d\beta\,d\theta\\
&=\lambda_{N-1}(B_j)
\int_{I\cap E_{j,m}}V_j(\theta)\,d\theta\\
&=(2R)^{N-1}
\int_{I\cap E_{j,m}}V_j(\theta)\,d\theta.
\end{aligned}
\label{eq:step-003-tonelli}
\end{equation}
Equations~\eqref{eq:step-003-localized-area} and
\eqref{eq:step-003-tonelli} prove \eqref{eq:step-003-area-bound}.  For
$N=1$, the inner integral is over the empty tuple with zero-dimensional
mass one.

For completeness, the multiplicity mechanism in the cited area formula is
explicit here.  The count
$N(\Psi_j,D_{j,m},\alpha)=\#\{x\in D_{j,m}:\Psi_j(x)=\alpha\}$ may be
finite or $+\infty$; it is at least one on the image, which is precisely the
direction used in \eqref{eq:step-003-localized-area}.  Distinct roots of one
coefficient vector are multiple preimages and only increase this count.  At
a represented root $\alpha=\Psi_j(\theta,\beta)$, write
$n=b+\sum_{i\neq j}\beta_iF_i$.  Since $T_j=-n/F_j$,
\begin{equation}
\partial_\theta T_j
=-\frac{n'+T_jF_j'}{F_j}
=-\frac{\phi_\alpha'(\theta)}{F_j(\theta)}.
\label{eq:step-003-critical-root}
\end{equation}
Thus a tangent or differentiably multiple root is a critical preimage.
Applying \eqref{eq:step-003-localized-area} to the measurable critical
subset $\{\partial_\theta T_j=0\}\cap D_{j,m}$ shows that its image has
zero $N$-volume.  If a coefficient has infinitely many roots, or its affine
combination vanishes on an interval, the multiplicity may be infinite;
the area formula remains valid in the extended nonnegative sense.  Neither
critical roots nor infinite fibers need be excluded.
\end{proof}

\begin{lemma}[Exact exhausted graph-image coverage]
\label{lem:step-003-exact-coverage}
Under Lemma~\ref{lem:step-001-root-feasible},
Proposition~\ref{prop:step-002-selector-partition},
Lemma~\ref{lem:step-002-chart-exhaustion},
Proposition~\ref{prop:step-002-exact-chart}, and
Lemma~\ref{lem:step-003-measurable-charts}, define
\begin{equation}
D_j
=\{(\theta,\beta):
\theta\in I\cap E_j,\ \beta\in B_j,\
|T_j(\theta,\beta)|\leq R\}.
\label{eq:step-003-full-domain}
\end{equation}
Then
\begin{equation}
D_{j,m}\uparrow D_j,
\qquad
\Psi_j(D_{j,m})\uparrow\Psi_j(D_j),
\label{eq:step-003-nested-images}
\end{equation}
and
\begin{equation}
S_I
=\bigcup_{j=1}^N\Psi_j(D_j)
=\bigcup_{j=1}^N\bigcup_{m=1}^{\infty}\Psi_j(D_{j,m}).
\label{eq:step-003-exact-coverage}
\end{equation}
The identities include roots at interval endpoints and coefficient-cube
boundaries, least-index pivot ties, and coefficients with one, multiple, or
infinitely many roots.  Empty $K_R$, empty cells, and $N=1$ are included.
\end{lemma}

\begin{proof}
Only the condition $\theta\in E_{j,m}$ in
\eqref{eq:step-003-exhausted-domain} depends on $m$.
Lemma~\ref{lem:step-002-chart-exhaustion} gives
$E_{j,m}\uparrow E_j$, while $I,B_j,T_j$, and $|T_j|\leq R$ remain fixed.
This proves the domain assertion in
\eqref{eq:step-003-nested-images}.  A fixed map sends an increasing union to
the increasing union of its images, proving the image assertion.

Take $\alpha\in S_I$ and choose any $\theta\in I$ with
$\phi_\alpha(\theta)=0$.  Lemma~\ref{lem:step-001-root-feasible} gives
$\theta\in K_R$.  Proposition~\ref{prop:step-002-selector-partition} puts
$\theta$ in one and only one least-index cell $E_j$ and gives
$F_j(\theta)\neq0$.  Set $\beta=\alpha_{-j}$ in increasing
original-coordinate order.  Then $\beta\in B_j$, and
Proposition~\ref{prop:step-002-exact-chart} gives
$T_j(\theta,\beta)=\alpha_j$.  Hence $(\theta,\beta)\in D_j$ and
$\Psi_j(\theta,\beta)=\alpha$.  The exhaustion lemma also puts this
$\theta$ in some finite $E_{j,m}$, so $\alpha$ belongs to an exhausted
image.

Conversely, if $\alpha=\Psi_j(\theta,\beta)$ with
$(\theta,\beta)\in D_j$, the beta coordinates lie in the closed cube and
$|T_j|\leq R$, so $\alpha\in[-R,R]^N$ coordinate by coordinate.  Also
$\theta\in I\cap E_j$, and the original-coordinate identity in
Proposition~\ref{prop:step-002-exact-chart} gives
$\phi_\alpha(\theta)=0$.  Thus $\alpha\in S_I$, proving
\eqref{eq:step-003-exact-coverage}.

The proof uses a root only as a witness, never as a unique or simple root.
Several or infinitely many roots yield several or infinitely many preimages
but one image point.  A pivot tie enters exactly one least-index cell.  All
cube inequalities are closed and the actual interval $I$ is used, so no
boundary point is discarded.  If $K_R=\varnothing$, root feasibility makes
$S_I$ and every chart set empty.  For $N=1$, the proof uses the unique empty
beta tuple and the scalar map $\Psi_1=T_1$.
\end{proof}

\begin{proposition}[Root-event coefficient-volume certificate]
\label{prop:step-003-volume-certificate}
Under Assumptions~\ref{assump:shared-pfaffian-chain} and
\ref{assump:no-forced-root}, Lemma~\ref{lem:step-001-root-feasible},
Proposition~\ref{prop:step-001-fixed-family-bound},
Proposition~\ref{prop:step-002-selector-partition},
Lemma~\ref{lem:step-002-chart-exhaustion},
Proposition~\ref{prop:step-003-area-bound}, and
Lemma~\ref{lem:step-003-exact-coverage}, every interval
$I\subseteq\Theta$ satisfies
\begin{equation}
\lambda_N(S_I)
\leq(2R)^{N-1}\Gamma_{\rm piv}(b,F;R)|I|.
\label{eq:step-003-volume-certificate}
\end{equation}
This includes empty $K_R$, empty cells, zero-length intervals, $N=1$,
cube and interval boundaries, tangent and multiple roots, and finite or
infinite root fibers.  There is no chart-count or multiplicity factor.
\end{proposition}

\begin{proof}
Fix $j$.  By Lemma~\ref{lem:step-003-exact-coverage}, the measurable sets
$\Psi_j(D_{j,m})$ increase to $\Psi_j(D_j)$.  Continuity from below,
Proposition~\ref{prop:step-003-area-bound}, and monotone convergence give
\begin{equation}
\begin{aligned}
\lambda_N(\Psi_j(D_j))
&=\lim_{m\to\infty}\lambda_N(\Psi_j(D_{j,m}))\\
&\leq(2R)^{N-1}\lim_{m\to\infty}
\int_{I\cap E_{j,m}}V_j(\theta)\,d\theta\\
&=(2R)^{N-1}\int_{I\cap E_j}V_j(\theta)\,d\theta.
\end{aligned}
\label{eq:step-003-exhausted-volume}
\end{equation}
Indeed, the functions that equal $V_j$ on $I\cap E_{j,m}$ and zero off
that set increase pointwise to the analogous function for $I\cap E_j$.
They are nonnegative and Borel; this formulation avoids any
$0\cdot(+\infty)$ convention.  No $m$-dependent constant survives.

Exact coverage, finite image subadditivity, and
\eqref{eq:step-003-exhausted-volume} yield
\begin{equation}
\lambda_N(S_I)
\leq(2R)^{N-1}
\sum_{j=1}^N\int_{I\cap E_j}V_j(\theta)\,d\theta.
\label{eq:step-003-chart-sum}
\end{equation}
On $E_j$, Proposition~\ref{prop:step-002-selector-partition} gives
$V_j=\min_kV_k\leq\Gamma_{\rm piv}(b,F;R)$.  Because the $E_j$ are
pairwise disjoint and have union $K_R$,
\begin{equation}
\begin{aligned}
\sum_{j=1}^N\int_{I\cap E_j}V_j(\theta)\,d\theta
&\leq\Gamma_{\rm piv}(b,F;R)\sum_{j=1}^N|I\cap E_j|\\
&=\Gamma_{\rm piv}(b,F;R)|I\cap K_R|\\
&\leq\Gamma_{\rm piv}(b,F;R)|I|.
\end{aligned}
\label{eq:step-003-disjoint-length}
\end{equation}
Combining \eqref{eq:step-003-chart-sum} and
\eqref{eq:step-003-disjoint-length} proves
\eqref{eq:step-003-volume-certificate}.  Thus the chart sum spends the
one-copy parameter-length budget $|I\cap K_R|$, not $N|I|$.

If $K_R=\varnothing$, every chart set is empty,
Lemma~\ref{lem:step-001-root-feasible} gives $\Gamma_{\rm piv}=0$, and both
sides of \eqref{eq:step-003-volume-certificate} vanish.  An empty cell
contributes zero.  If $|I|=0$,
\eqref{eq:step-003-disjoint-length} forces every chart integral to zero,
including for a singleton containing a root.  For $N=1$, there is one cell,
$(2R)^{N-1}=1$, and beta volume is one.  Multiplicity and critical-root
cases entered before the limits and chart sum, so no root count,
transversality, or finite-fiber condition appears.
\end{proof}

\begin{proof}[Completion of the coefficient-volume argument]
Lemma~\ref{lem:step-001-root-feasible} restricts every supported root to
$K_R$ and closes the empty branch.  Proposition
\ref{prop:step-002-selector-partition}, Lemma
\ref{lem:step-002-chart-exhaustion}, and Proposition
\ref{prop:step-002-exact-chart} supply the disjoint Borel cells, nonzero
selected pivots, finite-level exhaustion, original coefficient chart, and
velocity cap before the measure argument.

Lemma~\ref{lem:step-003-measurable-charts} proves event and chart
measurability and ambient local Lipschitz regularity, including endpoints.
Proposition~\ref{prop:step-003-area-bound} computes the signed determinant
in original coordinate order and applies the noninjective area formula on
disjoint measurable localizations.  Its beta integration gives exactly
$(2R)^{N-1}$, while its multiplicity and critical-set calculations cover
repeated roots, tangencies, differentiable multiple roots, and infinite
fibers.  Lemma~\ref{lem:step-003-exact-coverage} proves the graph-image
union identity, including every cube and interval boundary, and proves
nestedness in $m$.  Finally,
Proposition~\ref{prop:step-003-volume-certificate} combines continuity from
below, monotone convergence, and
\[
\sum_j|I\cap E_j|=|I\cap K_R|\leq|I|
\]
to obtain \eqref{eq:step-003-volume-certificate}, with no multiplicity,
boundary, localization, or chart-count factor.
\end{proof}

\subsection{Joint-Density Conversion and Uniform Ratio}
\label{app:step-004}

\begin{proposition}[Full joint-density conversion in coefficient space]
\label{prop:step-004-density-conversion}
Under Assumption~\ref{assump:joint-density-cap},
Lemma~\ref{lem:step-003-measurable-charts}, and
Proposition~\ref{prop:step-003-volume-certificate}, if
$\mu\in\mathcal D_{N,R,\kappa}$ and $I\subseteq\Theta$ is any interval with
$|I|>0$, then
\begin{equation}
\begin{aligned}
\Pr_{\alpha\sim\mu}\!\left[
  \exists\theta\in I:\phi_\alpha(\theta)=0
\right]
&=\int_{S_I}f_\mu(\alpha)\,d\lambda_N(\alpha)\\
&\leq
\kappa(2R)^{N-1}\Gamma_{\rm piv}(b,F;R)|I|\\
&=\frac{A\Gamma_{\rm piv}(b,F;R)}{2R}|I|.
\end{aligned}
\label{eq:step-004-density-conversion}
\end{equation}
This is an ordinary-probability statement for the original root event.  It
holds for arbitrarily correlated admissible laws, retains the original
coefficient dimension, and includes $N=1$, empty $K_R$, and every interval
endpoint convention covered by Proposition
\ref{prop:step-003-volume-certificate}.
\end{proposition}

\begin{proof}
Assumption~\ref{assump:joint-density-cap} gives
$f_\mu=0$ almost everywhere outside $[-R,R]^N$.  Hence
\begin{equation}
\mu([-R,R]^N)
=\int_{[-R,R]^N}f_\mu\,d\lambda_N=1.
\label{eq:step-004-cube-support}
\end{equation}
Therefore the unrestricted event in
\eqref{eq:step-004-density-conversion}, under $\alpha\sim\mu$, agrees almost
surely with its intersection with the coefficient cube, which is precisely
$S_I$.  Lemma~\ref{lem:step-003-measurable-charts} makes $S_I$ universally
measurable, so the probability and density integral are well-defined.

The pointwise full joint-density ceiling and
Proposition~\ref{prop:step-003-volume-certificate} now give, without
conditioning or coordinate factorization,
\begin{equation}
\begin{aligned}
\Pr_{\alpha\sim\mu}\!\left[
  \exists\theta\in I:\phi_\alpha(\theta)=0
\right]
&=\mu(S_I)\\
&=\int_{S_I}f_\mu(\alpha)\,d\lambda_N(\alpha)\\
&\leq\int_{S_I}\kappa\,d\lambda_N(\alpha)\\
&=\kappa\lambda_N(S_I)\\
&\leq\kappa(2R)^{N-1}
\Gamma_{\rm piv}(b,F;R)|I|.
\end{aligned}
\label{eq:step-004-density-chain}
\end{equation}
Because $R>0$ and $A=(2R)^N\kappa$, direct algebra gives
\begin{equation}
\kappa(2R)^{N-1}
=\frac{(2R)^N\kappa}{2R}
=\frac A{2R}.
\label{eq:step-004-a-identity}
\end{equation}
Equations~\eqref{eq:step-004-density-chain} and
\eqref{eq:step-004-a-identity} prove
\eqref{eq:step-004-density-conversion}.  Only the full joint density was
used, so arbitrary dependence among coordinates changes no line.

For $N=1$, the inherited beta-volume factor is $(2R)^0=1$, and
\eqref{eq:step-004-a-identity} becomes $A/(2R)=\kappa$.  Thus the bound
reads
\[
\Pr_{\alpha\sim\mu}\!\left[
  \exists\theta\in I:\phi_\alpha(\theta)=0
\right]
\leq\kappa\Gamma_{\rm piv}(b,F;R)|I|
\]
with no exceptional convention or hidden factor.  If $K_R=\varnothing$,
Lemma~\ref{lem:step-001-root-feasible} gives $S_I=\varnothing$ and
$\Gamma_{\rm piv}=0$, so every line is zero.  Open, closed, half-open, and
relative-endpoint interval conventions are unchanged because
\eqref{eq:step-004-density-chain} integrates the same $S_I$, not an
enlargement or closure.  A right-hand side larger than one remains a valid
probability upper bound; no small-interval threshold or truncation is
introduced.
\end{proof}

\begin{proposition}[Uniform finite anti-concentration ratio]
\label{prop:step-004-uniform-ratio}
Under Assumption~\ref{assump:joint-density-cap},
Proposition~\ref{prop:step-001-fixed-family-bound}, and
Proposition~\ref{prop:step-004-density-conversion},
\begin{equation}
\sup_{\mu\in\mathcal D_{N,R,\kappa}}
\sup_{\substack{I\subseteq\Theta\text{ an interval}\\|I|>0}}
\frac{\Pr_{\alpha\sim\mu}[
  \exists\theta\in I:\phi_\alpha(\theta)=0]}{|I|}
\leq
\frac{A\Gamma_{\rm piv}(b,F;R)}{2R}
<\infty.
\label{eq:step-004-uniform-ratio}
\end{equation}
The conclusion is uniform over the full, possibly correlated law class and
over every positive-length interval, with no hidden constant.
\end{proposition}

\begin{proof}
For every admissible $\mu$ and every stated interval $I$, one has $|I|>0$.
Dividing the last inequality in
Proposition~\ref{prop:step-004-density-conversion} by this positive number
gives
\begin{equation}
\frac{\Pr_{\alpha\sim\mu}[
  \exists\theta\in I:\phi_\alpha(\theta)=0]}{|I|}
\leq
\frac{A\Gamma_{\rm piv}(b,F;R)}{2R}.
\label{eq:step-004-pairwise-ratio}
\end{equation}
The right-hand side is independent of both $\mu$ and $I$.  It therefore
bounds the inner set of ratios for each law and then the outer set of inner
suprema.  Taking the two least upper bounds in that order proves
\eqref{eq:step-004-uniform-ratio}.  If either indexing class is empty, no
admissible pair can violate the same uniform upper-bound statement.

Finally, $A=(2R)^N\kappa<\infty$, $R>0$, and
Proposition~\ref{prop:step-001-fixed-family-bound}, together with the empty
branch in Lemma~\ref{lem:step-001-root-feasible}, gives
$\Gamma_{\rm piv}(b,F;R)<\infty$.  Hence the right-hand side is finite.  No
union bound, conditioning, expectation, independence reduction,
marginalization, or probability-mode conversion occurs when taking the
suprema.
\end{proof}

\begin{proof}[Completion of the density-conversion argument]
Lemma~\ref{lem:step-003-measurable-charts} supplies measurability of the
cube-restricted root event, and
Proposition~\ref{prop:step-003-volume-certificate} supplies
\[
\lambda_N(S_I)
\leq(2R)^{N-1}\Gamma_{\rm piv}(b,F;R)|I|
\]
in the original coefficient dimension, including all boundary,
multiplicity, and $N=1$ cases.

Proposition~\ref{prop:step-004-density-conversion} applies the primitive
full joint-density cap to this identical event and proves, for every
admissible law and every positive-length interval,
\[
\Pr_{\alpha\sim\mu}(S_I)
\leq\kappa(2R)^{N-1}\Gamma_{\rm piv}(b,F;R)|I|
=\frac{A\Gamma_{\rm piv}(b,F;R)}{2R}|I|.
\]
Because the proof uses neither marginal nor conditional laws, arbitrary
correlation is preserved.  Proposition~\ref{prop:step-004-uniform-ratio}
then divides only by $|I|>0$ and takes the interval and law suprema without
changing the constant.  Fixed-family finiteness of $\Gamma_{\rm piv}$ makes
the resulting anti-concentration bound finite.
\end{proof}

\subsection{Exact Scale-Stress Calculation}
\label{app:step-005}

\begin{proposition}[Exact scale-stress pivot profile]
\label{prop:step-005-pivot-profile}
Under Assumptions~\ref{assump:shared-pfaffian-chain} and
\ref{assump:no-forced-root}, Lemma~\ref{lem:step-001-root-feasible}, and
Proposition~\ref{prop:step-001-fixed-family-bound}, suppose
\[
0<\delta\leq1,qquad
\Theta=[-1,1],qquad b(\theta)=0,qquad
F(\theta)=(1,\theta/\delta),qquad R=1.
\]
Then $K_1=[-1,1]$ and
\begin{equation}
V_1(\theta)=\frac1\delta
\quad(-1\leq\theta\leq1),
\label{eq:step-005-v1}
\end{equation}
while
\begin{equation}
V_2(0)=+\infty,
\qquad
V_2(\theta)=\frac\delta{\theta^2}
\quad(\theta\neq0).
\label{eq:step-005-v2}
\end{equation}
At $\theta=0$, coordinate $1$ is the unique finite minimizing pivot.  For
$\theta\neq0$, coordinate $1$ uniquely minimizes when
$0<|\theta|<\delta$, both coordinates minimize when
$|\theta|=\delta$, and coordinate $2$ uniquely minimizes when
$\delta<|\theta|\leq1$.  Under the least-index tie rule,
$j_*(\theta)=1$ at $|\theta|=\delta$.  Consequently,
\begin{equation}
\Gamma_{\rm piv}(0,(1,\theta/\delta);1)=\frac1\delta.
\label{eq:step-005-exact-gamma}
\end{equation}
\end{proposition}

\begin{proof}
For every $\theta\in[-1,1]$,
\[
|b(\theta)|=0
\leq |F_1(\theta)|+|F_2(\theta)|
=1+\frac{|\theta|}{\delta}.
\]
Thus the definition gives $K_1=\Theta=[-1,1]$, consistently with
Lemma~\ref{lem:step-001-root-feasible}.

For the constant pivot $F_1=1$, the definition of $V_1$ and $R=1$ give
\[
\begin{aligned}
V_1(\theta)
&=\left|\left(\frac01\right)'\right|
 +\left|\left(\frac{\theta/\delta}{1}\right)'\right|
=\frac1\delta.
\end{aligned}
\]
The second pivot vanishes exactly at $\theta=0$, so the extended definition
gives $V_2(0)=+\infty$.  On its ordinary domain $\theta\neq0$,
\[
\begin{aligned}
V_2(\theta)
&=\left|\left(\frac0{\theta/\delta}\right)'\right|
 +\left|\left(\frac1{\theta/\delta}\right)'\right|\\
&=0+\left|\left(\frac\delta\theta\right)'\right|
=\frac\delta{\theta^2}.
\end{aligned}
\]
This also shows directly that coordinate $1$ is the unique finite minimizer
at zero.  Away from zero,
\[
\frac\delta{\theta^2}
\begin{cases}
>1/\delta,&0<|\theta|<\delta,\\
=1/\delta,&|\theta|=\delta,\\
<1/\delta,&\delta<|\theta|\leq1.
\end{cases}
\]
Therefore
\begin{equation}
\min\{V_1(\theta),V_2(\theta)\}
=
\begin{cases}
1/\delta,&|\theta|\leq\delta,\\
\delta/\theta^2,&\delta<|\theta|\leq1.
\end{cases}
\label{eq:step-005-min-profile}
\end{equation}
When $\delta=1$, the second branch is empty and the pivots tie only at
$\theta=\pm1$; \eqref{eq:step-005-min-profile} and the least-index rule
remain valid.  The pointwise minimum never exceeds $1/\delta$ and equals
$1/\delta$ throughout $[-\delta,\delta]$, in particular at zero.  Taking
the supremum over $K_1=[-1,1]$ proves
\eqref{eq:step-005-exact-gamma}.
\end{proof}

\begin{lemma}[Exact closed-wedge characterization of the root event]
\label{lem:step-005-wedge-event}
Under Assumptions~\ref{assump:shared-pfaffian-chain} and
\ref{assump:no-forced-root}, suppose
$\Theta=[-1,1]$, $b(\theta)=0$,
$F(\theta)=(1,\theta/\delta)$, and
$0<\epsilon\leq\delta\leq1$.  Then, as an equality of subsets of
$[-1,1]^2$,
\begin{equation}
\begin{aligned}
&\left\{\alpha:\exists\theta\in[0,\epsilon],\
\alpha_1+\alpha_2\theta/\delta=0\right\}\\
&\qquad=
\left\{\alpha:\alpha_1\alpha_2\leq0,\
|\alpha_1|\leq\frac\epsilon\delta|\alpha_2|\right\}.
\end{aligned}
\label{eq:step-005-wedge-event}
\end{equation}
The equality includes all coefficient-axis cases.  The non-strict sign and
magnitude inequalities correspond, respectively, to the closed interval
endpoints $\theta=0$ and $\theta=\epsilon$.
\end{lemma}

\begin{proof}
Put $r=\epsilon/\delta$, so $0<r\leq1$.  Since $\delta>0$, the change of
variable $t=\theta/\delta$ maps $[0,\epsilon]$ bijectively onto $[0,r]$.
Thus the event on the left of \eqref{eq:step-005-wedge-event} is exactly
\[
\{\alpha:\exists t\in[0,r],\ \alpha_1+\alpha_2t=0\}.
\]

Suppose first that $\alpha_2\neq0$.  The only possible root parameter is
$t=-\alpha_1/\alpha_2$, and
\[
-\frac{\alpha_1}{\alpha_2}\in[0,r]
\quad\Longleftrightarrow\quad
\alpha_1\alpha_2\leq0
\quad\text{and}\quad
|\alpha_1|\leq r|\alpha_2|.
\]
The weak inequality at zero includes $t=0$, equivalently $\theta=0$, and
equality in the magnitude constraint includes $t=r$, equivalently
$\theta=\epsilon$.

Now suppose $\alpha_2=0$.  A root exists if and only if $\alpha_1=0$.
On the proposed wedge, the magnitude inequality becomes
$|\alpha_1|\leq0$, with the same conclusion.  On the other coefficient
axis, $\alpha_1=0$ gives a root at $\theta=0$ for every $\alpha_2$, and
both wedge inequalities hold.  At the origin, every
$\theta\in[0,\epsilon]$ is a root.  The set identity is therefore exact on
the axes as well as off them; no null-boundary modification is used.
\end{proof}

\begin{proposition}[Exact planar wedge probability]
\label{prop:step-005-wedge-probability}
Under Assumption~\ref{assump:joint-density-cap} specialized to the uniform
law on $[-1,1]^2$, Lemma~\ref{lem:step-005-wedge-event}, and
$0<\epsilon\leq\delta$, the coefficient set having a root in
$[0,\epsilon]$ has two-dimensional Lebesgue measure $\epsilon/\delta$.
Consequently,
\begin{equation}
\Pr\!\left[\exists\theta\in[0,\epsilon]:
\phi_\alpha(\theta)=0\right]
=\frac\epsilon{4\delta}.
\label{eq:step-005-wedge-probability}
\end{equation}
\end{proposition}

\begin{proof}
Again set $r=\epsilon/\delta\in(0,1]$.  By
Lemma~\ref{lem:step-005-wedge-event}, the root event is the closed set
\[
W_r=\{\alpha\in[-1,1]^2:
\alpha_1\alpha_2\leq0,\ |\alpha_1|\leq r|\alpha_2|\}.
\]
It is the union of the two closed triangular components
\[
W_r^+=\{0\leq\alpha_2\leq1,
-r\alpha_2\leq\alpha_1\leq0\},
\]
\[
W_r^-=\{-1\leq\alpha_2\leq0,
0\leq\alpha_1\leq-r\alpha_2\}.
\]
The condition $r\leq1$ ensures that every displayed horizontal endpoint
lies in $[-1,1]$, so neither triangle is clipped by a vertical side of the
coefficient square.  Their intersection is only the origin.  Tonelli's
theorem applied to the nonnegative indicators and the displayed vertical
sections gives
\begin{equation}
\lambda_2(W_r^+)
=\int_0^1r\alpha_2\,d\alpha_2
=\frac r2,
\label{eq:step-005-upper-wedge}
\end{equation}
and
\begin{equation}
\lambda_2(W_r^-)
=\int_{-1}^0(-r\alpha_2)\,d\alpha_2
=\frac r2.
\label{eq:step-005-lower-wedge}
\end{equation}
Their intersection has area zero, and therefore
\begin{equation}
\lambda_2(W_r)=r=\frac\epsilon\delta.
\label{eq:step-005-wedge-area}
\end{equation}

For completeness, the coefficient axes are one-dimensional subsets of the
plane and have zero planar measure.  Under the vertical-section convention
$B_y=\{x:(x,y)\in B\}$, the vertical axis has singleton sections for every
$y$, whereas the horizontal axis has section $[-1,1]$ only at $y=0$ and
empty sections otherwise.  Its sole nonzero section occurs on a null set of
slice parameters, so its area is zero.  The sloping endpoint edges also have
singleton vertical sections and area zero.  This prevents boundary double
counting in \eqref{eq:step-005-upper-wedge}--
\eqref{eq:step-005-lower-wedge}, while
Lemma~\ref{lem:step-005-wedge-event} already established that the axes and
edges belong to the event with the correct closed-endpoint convention.

Finally, the uniform law has density $1/4$ on the square.  Hence
\[
\begin{aligned}
\Pr\!\left[\exists\theta\in[0,\epsilon]:
\phi_\alpha(\theta)=0\right]
&=\int_{W_r}\frac14\,d\alpha\\
&=\frac14\lambda_2(W_r)
=\frac\epsilon{4\delta}.
\end{aligned}
\]
There is no conditioning, union bound, or independence step in this
equality.
\end{proof}

\begin{proof}[Completion of the scale-stress calculation]
Lemma~\ref{lem:step-001-root-feasible} and
Proposition~\ref{prop:step-001-fixed-family-bound} place the scale-stress
specialization inside the compact root-feasible, finite-conditioning
regime.  Proposition~\ref{prop:step-005-pivot-profile} sharpens this to an
exact calculation: the constant pivot has speed $1/\delta$, the second
pivot has speed $\delta/\theta^2$ away from zero and is unavailable at zero,
and their pointwise minimum has supremum exactly $1/\delta$.  The piecewise
description covers zero, both regions away from zero, the tie
$|\theta|=\delta$, and the endpoint case $\delta=1$.

For the probability statement, Lemma~\ref{lem:step-005-wedge-event}
identifies the original closed-interval event exactly with the two
opposite-sign wedges.  It accounts for roots at $\theta=0$ and
$\theta=\epsilon$, the full $\alpha_1=0$ axis segment, and the fact that the
$\alpha_2=0$ axis contributes only the origin.  Proposition
\ref{prop:step-005-wedge-probability} computes the two wedge areas as
$\epsilon/(2\delta)$ each and multiplies their total by the uniform density
$1/4$.  Thus the fixed scale-stress family has
$\Gamma_{\rm piv}=1/\delta$ and root probability
$\epsilon/(4\delta)$, both with their literal constants.
\end{proof}

\subsection{Exact Monic Lower-Coefficient Specialization}
\label{app:step-006}

\begin{proposition}[Exact lower-coefficient affine representation]
\label{prop:step-006-family}
Let $d\in\mathbb Z$, $d\geq1$, $R>0$, and
$\kappa\in(0,\infty)$, and impose
Assumption~\ref{assump:joint-density-cap} on the $d$-dimensional law of
$\alpha=(\alpha_0,\ldots,\alpha_{d-1})$.  On any compact interval
$\Theta\subset\mathbb R$, define
\[
b_d(\theta)=\theta^d,
\qquad
F_d(\theta)=(1,\theta,\ldots,\theta^{d-1}).
\]
Then, for every $\theta\in\Theta$ and $\alpha\in\mathbb R^d$,
\begin{equation}
b_d(\theta)+\langle\alpha,F_d(\theta)\rangle
=\theta^d+\sum_{k=0}^{d-1}\alpha_k\theta^k
=p_\alpha(\theta).
\label{eq:step-006-object-identity}
\end{equation}
This family satisfies Assumptions~\ref{assump:shared-pfaffian-chain} and
\ref{assump:no-forced-root}, and its descriptors are exactly
\begin{equation}
q=0,
\qquad M=0,
\qquad N=d,
\qquad \Delta_{\rm rnd}=d-1,
\qquad \Delta_{\rm aff}=d,
\qquad A=(2R)^d\kappa.
\label{eq:step-006-descriptors}
\end{equation}
Only the $d$ lower coefficients are random.
\end{proposition}

\begin{proof}
The coordinate pairing is $F_{k+1}(\theta)=\theta^k$ for
$0\leq k\leq d-1$.  Therefore
\[
\langle\alpha,F_d(\theta)\rangle
=\sum_{k=0}^{d-1}\alpha_k\theta^k,
\]
which proves every equality in \eqref{eq:step-006-object-identity}.  The
coefficient of $\theta^d$ is the fixed value one inside $b_d$; it is not a
coordinate of $\alpha$.

All displayed functions are polynomials and hence $C^1$ on every compact
$\Theta$.  They have a common Pfaffian representation with no chain: take
$q=0$, so the convention gives $M=0$, and use the ordinary polynomials
$Q_0(\theta)=\theta^d$ and
$Q_{k+1}(\theta)=\theta^k$.  This verifies
Assumption~\ref{assump:shared-pfaffian-chain}.  There are $d$ random
outputs, so $N=d$.  Their degrees are $0,1,\ldots,d-1$, giving
$\Delta_{\rm rnd}=d-1$.  Including the deterministic offset of degree $d$
gives $\Delta_{\rm aff}=d$.

The first coordinate of $F_d(\theta)$ equals one for every $\theta$.
Consequently $F_d(\theta)\neq0$, and hence
$(b_d(\theta),F_d(\theta))\neq(0,0)$ pointwise.  This verifies
Assumption~\ref{assump:no-forced-root}, including at $\theta=0$.  Finally,
substituting $N=d$ into $A=(2R)^N\kappa$ gives the last identity in
\eqref{eq:step-006-descriptors}.  None of these identifications changes the
law: Assumption~\ref{assump:joint-density-cap} remains a full joint density
on $\mathbb R^d$, with arbitrary correlation among the lower coefficients.
\end{proof}

\begin{lemma}[Constant and top monomial pivot certificate]
\label{lem:step-006-two-pivots}
Under Proposition~\ref{prop:step-006-family}, the following conclusions hold
for every $\theta\in\mathbb R$.  If $d=1$, the sole feature is constant and
\begin{equation}
V_{\rm const}(\theta)=1.
\label{eq:step-006-d1-pivot}
\end{equation}
If $d\geq2$, then on $|\theta|\leq1$ the constant feature is a legal pivot
and
\begin{equation}
V_{\rm const}(\theta)
=d|\theta|^{d-1}
+R\sum_{k=1}^{d-1}k|\theta|^{k-1}
\leq d+\frac{Rd(d-1)}2,
\label{eq:step-006-constant-pivot}
\end{equation}
whereas on $|\theta|>1$ the top feature $\theta^{d-1}$ is a legal pivot and
\begin{equation}
V_{\rm top}(\theta)
=1+R\sum_{m=1}^{d-1}\frac{m}{|\theta|^{m+1}}
\leq1+\frac{Rd(d-1)}2
\leq d+\frac{Rd(d-1)}2.
\label{eq:step-006-top-pivot}
\end{equation}
Consequently, for every $d\geq1$ and every $\theta\in\mathbb R$,
\begin{equation}
\min_{1\leq j\leq d}V_j(\theta)
\leq d+\frac{Rd(d-1)}2.
\label{eq:step-006-pointwise-min}
\end{equation}
The two regions are disjoint, their common boundary is assigned only to the
constant pivot, and the top pivot is never used near its zero.
\end{lemma}

\begin{proof}
First suppose $d=1$.  Then $b_1(\theta)=\theta$ and $F_1(\theta)=1$.
There are no nonpivot features, so
\[
V_{\rm const}(\theta)
=\left|\left(\frac\theta1\right)'\right|=1.
\]
This proves \eqref{eq:step-006-d1-pivot}, including $\theta=0$, and its
right-hand side equals $d+Rd(d-1)/2=1$.

Now suppose $d\geq2$.  The constant feature is nonzero on all of
$\mathbb R$.  Pivoting on it gives the exact ratios
\[
\frac{b_d(\theta)}1=\theta^d,
\qquad
\frac{\theta^k}{1}=\theta^k
\quad(1\leq k\leq d-1).
\]
Direct differentiation gives the equality in
\eqref{eq:step-006-constant-pivot}.  When $|\theta|\leq1$,
\[
d|\theta|^{d-1}\leq d,
\qquad
k|\theta|^{k-1}\leq k
\quad(1\leq k\leq d-1),
\]
where the $k=1$ derivative equals one also at $\theta=0$.  Hence
\begin{equation}
V_{\rm const}(\theta)
\leq d+R\sum_{k=1}^{d-1}k
=d+\frac{Rd(d-1)}2.
\label{eq:step-006-constant-sum}
\end{equation}

For $|\theta|>1$, the top feature $\theta^{d-1}$ is nonzero.  The offset
ratio is
\[
\frac{\theta^d}{\theta^{d-1}}=\theta,
\]
whose derivative has magnitude one.  For every other feature $\theta^k$,
$0\leq k\leq d-2$, set $m=d-1-k$.  Then $1\leq m\leq d-1$ and
\[
\frac{\theta^k}{\theta^{d-1}}=\theta^{-m},
\qquad
| (\theta^{-m})' |
=\frac{m}{|\theta|^{m+1}}.
\]
Reindexing the feature sum proves the equality in
\eqref{eq:step-006-top-pivot}.  Since $|\theta|>1$, every denominator in
that sum is at least one, so
\begin{equation}
V_{\rm top}(\theta)
\leq1+R\sum_{m=1}^{d-1}m
=1+\frac{Rd(d-1)}2
\leq d+\frac{Rd(d-1)}2,
\label{eq:step-006-top-sum}
\end{equation}
where the last inequality uses $d\geq2$.

At $|\theta|=1$, \eqref{eq:step-006-constant-sum} applies and the top
region does not.  At $\theta=0$, only the everywhere-nonzero constant pivot
is used.  Thus the singularity of the quotient by $\theta^{d-1}$ at zero
never enters.  The conditions $|\theta|\leq1$ and $|\theta|>1$ form a
disjoint exhaustive partition, so no point or length is counted twice.  On
each region, the minimum over all legal pivots is at most the selected-pivot
value.  Equations~\eqref{eq:step-006-d1-pivot},
\eqref{eq:step-006-constant-sum}, and \eqref{eq:step-006-top-sum} prove
\eqref{eq:step-006-pointwise-min}.
\end{proof}

\begin{proposition}[Compact-independent monic conditioning bound]
\label{prop:step-006-compact-gamma}
Under Assumptions~\ref{assump:shared-pfaffian-chain} and
\ref{assump:no-forced-root}, as verified by
Proposition~\ref{prop:step-006-family}, Lemma
\ref{lem:step-001-root-feasible}, and
Lemma~\ref{lem:step-006-two-pivots}, let
$\Theta\subset\mathbb R$ be any compact interval and form $K_R$ and
$\Gamma_{\rm piv}(b_d,F_d;R)$ from the restrictions of $b_d,F_d$ to
$\Theta$.  Then
\begin{equation}
\Gamma_{\rm piv}(b_d,F_d;R)
\leq d+\frac{Rd(d-1)}2.
\label{eq:step-006-compact-gamma}
\end{equation}
The right-hand side is independent of $\Theta$.  Moreover, for every
interval $J\subseteq\Theta$ with $J\cap K_R=\varnothing$, no
$\alpha\in[-R,R]^d$ has a root of $p_\alpha$ in $J$.
\end{proposition}

\begin{proof}
Lemma~\ref{lem:step-001-root-feasible} applies because
Proposition~\ref{prop:step-006-family} verifies its regularity premise.  If
$K_R=\varnothing$, the definition gives
$\Gamma_{\rm piv}(b_d,F_d;R)=0$, so
\eqref{eq:step-006-compact-gamma} follows because its right-hand side is
positive.

Suppose instead that $K_R\neq\varnothing$.  For every $\theta\in K_R$,
Lemma~\ref{lem:step-006-two-pivots} gives
\[
\min_{1\leq j\leq d}V_j(\theta)
\leq d+\frac{Rd(d-1)}2.
\]
Taking the supremum over $\theta\in K_R$ preserves the same constant and
proves \eqref{eq:step-006-compact-gamma}.  The proof used only the value of
$\theta$, the split at $|\theta|=1$, and $d,R$; no endpoint, diameter, or
other quantity associated with $\Theta$ appears.  Thus the bound is uniform
over the auxiliary compact localization.

Finally, Lemma~\ref{lem:step-001-root-feasible} says that every root produced
by $\alpha\in[-R,R]^d$ must lie in $K_R$.  Therefore an interval $J$
disjoint from $K_R$ contains no such root.  This includes the empty-$K_R$
regime.
\end{proof}

\begin{lemma}[Zero-length bounded intervals are null]
\label{lem:step-006-zero-length}
Under Assumption~\ref{assump:joint-density-cap} on the $d$-dimensional
lower-coefficient vector and Proposition~\ref{prop:step-006-family}, if a
bounded interval $I\subset\mathbb R$ has $|I|=0$, then
\begin{equation}
\Pr_{\alpha\sim\mu}
[\exists\theta\in I:p_\alpha(\theta)=0]=0.
\label{eq:step-006-zero-length}
\end{equation}
If $I$ is nonempty, its root event is a proper affine hyperplane in the
original $\mathbb R^d$ coefficient space because the coefficient of
$\alpha_0$ is the constant feature one.
\end{lemma}

\begin{proof}
A nonempty interval containing two distinct points contains the segment
between them and therefore has positive length.  Hence a zero-length
interval is either empty or a singleton.

If $I=\varnothing$, the event in \eqref{eq:step-006-zero-length} is empty.
Otherwise write $I=\{\theta_0\}$.  By
\eqref{eq:step-006-object-identity}, the root event is
\begin{equation}
H_{\theta_0}
:=\left\{\alpha\in\mathbb R^d:
\alpha_0+\sum_{k=1}^{d-1}\alpha_k\theta_0^k=-\theta_0^d
\right\}.
\label{eq:step-006-hyperplane}
\end{equation}
The coefficient of $\alpha_0$ is one, so the defining affine functional is
nonconstant and $H_{\theta_0}$ is proper.

For $d=1$, \eqref{eq:step-006-hyperplane} is the singleton
$\{\alpha_0=-\theta_0\}$, which has one-dimensional Lebesgue measure zero.
For $d\geq2$, write
$\beta=(\alpha_1,\ldots,\alpha_{d-1})$.  Then
\eqref{eq:step-006-hyperplane} is the graph
\[
\alpha_0
=-\theta_0^d-\sum_{k=1}^{d-1}\beta_k\theta_0^k.
\]
This graph is closed, and Tonelli's theorem applied to its indicator gives
\begin{equation}
\begin{aligned}
\lambda_d(H_{\theta_0})
&=\int_{\mathbb R^{d-1}}
\lambda_1\!\left(
\left\{-\theta_0^d-
\sum_{k=1}^{d-1}\beta_k\theta_0^k\right\}
\right)d\beta\\
&=0.
\end{aligned}
\label{eq:step-006-hyperplane-nullity}
\end{equation}
Assumption~\ref{assump:joint-density-cap} supplies a full joint density
$f_\mu$ on this same $\mathbb R^d$, so
\[
\Pr_{\alpha\sim\mu}(H_{\theta_0})
=\int_{H_{\theta_0}}f_\mu(\alpha)\,d\lambda_d(\alpha)=0.
\]
This calculation neither factors the density nor adds the deterministic
leading coefficient.  It proves \eqref{eq:step-006-zero-length} for
arbitrary admissible correlation and also when the hyperplane meets a face
of the support cube.
\end{proof}

\begin{proposition}[Exact affine-monic probability recovery]
\label{prop:step-006-monic-bound}
Under Assumption~\ref{assump:joint-density-cap},
Proposition~\ref{prop:step-006-family},
Proposition~\ref{prop:step-006-compact-gamma},
Lemma~\ref{lem:step-006-zero-length}, and
Proposition~\ref{prop:step-004-density-conversion}, for every integer
$d\geq1$, every $\mu\in\mathcal D_{d,R,\kappa}$, and every bounded interval
$I\subset\mathbb R$,
\begin{equation}
\Pr_{\alpha\sim\mu}
[\exists\theta\in I:p_\alpha(\theta)=0]
\leq
\kappa(2R)^{d-1}
\left(d+\frac{Rd(d-1)}2\right)|I|.
\label{eq:step-006-monic-bound}
\end{equation}
This is ordinary probability for the $d$-dimensional, possibly correlated
lower-coefficient law and contains no hidden constant or auxiliary
condition.
\end{proposition}

\begin{proof}
If $|I|=0$, Lemma~\ref{lem:step-006-zero-length} makes the left-hand side of
\eqref{eq:step-006-monic-bound} zero, and its right-hand side is also zero.
Thus the inequality holds for both the empty and singleton cases.

Suppose $|I|>0$.  Because $I$ is bounded, its closure is a compact interval.
Choose any compact interval $\Theta$ containing that closure and restrict
$b_d,F_d$ to $\Theta$.  This choice introduces no rate parameter.
Proposition~\ref{prop:step-006-family} verifies the deterministic
assumptions of the affine theorem on $\Theta$, and
Assumption~\ref{assump:joint-density-cap} is already the required law on
$\mathbb R^d$.  Proposition~\ref{prop:step-006-compact-gamma} gives
\begin{equation}
\Gamma_{\rm piv}(b_d,F_d;R)
\leq d+\frac{Rd(d-1)}2
\label{eq:step-006-localized-gamma}
\end{equation}
for this localization, independently of its endpoints.

Apply Proposition~\ref{prop:step-004-density-conversion} with $N=d$.  By
the object equality \eqref{eq:step-006-object-identity}, its event is the
event in \eqref{eq:step-006-monic-bound}, not an approximation or
transformed event.  Hence
\begin{equation}
\begin{aligned}
\Pr_{\alpha\sim\mu}
[\exists\theta\in I:p_\alpha(\theta)=0]
&=\Pr_{\alpha\sim\mu}
[\exists\theta\in I:b_d(\theta)
  +\langle\alpha,F_d(\theta)\rangle=0]\\
&\leq\kappa(2R)^{d-1}
  \Gamma_{\rm piv}(b_d,F_d;R)|I|\\
&\leq\kappa(2R)^{d-1}
  \left(d+\frac{Rd(d-1)}2\right)|I|.
\end{aligned}
\label{eq:step-006-literal-substitution}
\end{equation}
The first inequality is the general affine density conversion; the second is
\eqref{eq:step-006-localized-gamma}.  Equivalently, its other constant form
reduces by direct algebra to the same factor:
\begin{equation}
\frac A{2R}
=\frac{(2R)^d\kappa}{2R}
=\kappa(2R)^{d-1}.
\label{eq:step-006-a-substitution}
\end{equation}

If $I\cap K_R=\varnothing$ for the chosen localization,
Lemma~\ref{lem:step-001-root-feasible} and
Proposition~\ref{prop:step-006-compact-gamma} show directly that the
cube-restricted event in \eqref{eq:step-006-literal-substitution} is empty.
Assumption~\ref{assump:joint-density-cap} assigns probability one to that
cube, so the unrestricted sampled event has probability zero.  Thus
\eqref{eq:step-006-monic-bound} also covers intervals wholly outside the
root-feasible region.  If $K_R=\varnothing$, the same argument gives
probability zero and $\Gamma_{\rm piv}=0$.  The cases $d=1$, $\theta=0$,
and $|\theta|=1$ are already included in
Lemma~\ref{lem:step-006-two-pivots}; no limiting or exceptional case
remains.
\end{proof}

\begin{proof}[Completion of the monic specialization]
Proposition~\ref{prop:step-006-family} identifies the setting objects
pointwise through
\[
b_d+\langle\alpha,F_d\rangle=p_\alpha
\]
with the deterministic monic coefficient outside the random vector, and it
establishes the descriptor tuple
\[
(q,M,N,\Delta_{\rm rnd},\Delta_{\rm aff},A)
=(0,0,d,d-1,d,(2R)^d\kappa).
\]
Thus the affine and monic events live in the same $d$-dimensional coefficient
probability space.

Lemma~\ref{lem:step-006-two-pivots} gives the constant-pivot formula on
$|\theta|\leq1$, the top-pivot formula on $|\theta|>1$, and the separate
$d=1$ identity $V_{\rm const}=1$.  The constant region includes zero and
$|\theta|=1$, while the top pivot is used only where
$\theta^{d-1}\neq0$.  Because the regions are disjoint and
$\Gamma_{\rm piv}$ takes a pointwise pivot minimum, the split incurs no
duplicate boundary or chart-count cost.  Proposition
\ref{prop:step-006-compact-gamma} then yields on every compact $\Theta$
\[
\Gamma_{\rm piv}(b_d,F_d;R)
\leq d+\frac{Rd(d-1)}2,
\]
independently of that compact interval and with empty or missed feasible
regions handled by root feasibility.

For positive-length bounded intervals,
Proposition~\ref{prop:step-006-monic-bound} localizes to a compact interval,
applies Proposition~\ref{prop:step-004-density-conversion}, and substitutes
the preceding conditioning certificate.  Lemma
\ref{lem:step-006-zero-length} supplies the empty and singleton cases.
Together they prove for every bounded interval
\[
\Pr_{\alpha\sim\mu}
[\exists\theta\in I:p_\alpha(\theta)=0]
\leq
\kappa(2R)^{d-1}
\left(d+\frac{Rd(d-1)}2\right)|I|,
\]
with arbitrary lower-coefficient correlation, ordinary probability, and no
leading-coordinate augmentation or conservative remainder.
\end{proof}

\subsection{Proof of the Main Theorem}
\label{app:main-proof}

\begin{proof}[Proof of Theorem~\ref{thm:main}]
Lemma~\ref{lem:step-001-root-feasible}, Lemma
\ref{lem:step-001-pivot-margin}, and
Proposition~\ref{prop:step-001-fixed-family-bound} prove the finiteness of
$\Gamma_{\rm piv}(b,F;R)$ in both the empty and nonempty feasible-set
branches.  Proposition~\ref{prop:step-004-density-conversion} gives the two
equal forms of the general probability bound for every admissible law and
positive-length interval, and
Proposition~\ref{prop:step-004-uniform-ratio} gives the nested uniform ratio.
This proves part~\textnormal{(i)}.

Proposition~\ref{prop:step-005-pivot-profile} proves
$\Gamma_{\rm piv}(b_\delta,F_\delta;1)=1/\delta$, and
Proposition~\ref{prop:step-005-wedge-probability} proves the exact uniform-law
probability $\epsilon/(4\delta)$ for
$0<\epsilon\leq\delta\leq1$.  This proves part~\textnormal{(ii)}.

Proposition~\ref{prop:step-006-family} verifies the exact monic object and
keeps its deterministic leading coefficient outside the $d$-dimensional
random law.  Proposition~\ref{prop:step-006-compact-gamma} proves the stated
conditioning bound on every compact localization, and
Proposition~\ref{prop:step-006-monic-bound} proves the claimed probability
inequality for every bounded interval, including zero-length intervals.
This proves part~\textnormal{(iii)} and completes the proof.
\end{proof}
